
\documentclass{article}
\usepackage{colortbl}
\usepackage{makecell}
\usepackage{multirow}
\usepackage{supertabular}

\begin{document}

\newcounter{utterance}

\centering \large Interaction Transcript for game `cladder', experiment `full\_v1.5\_default', episode 916 with selfrefine{-}Qwen3.5{-}27B.
\vspace{24pt}

{ \footnotesize  \setcounter{utterance}{1}
\setlength{\tabcolsep}{0pt}
\begin{supertabular}{c@{$\;$}|p{.15\linewidth}@{}p{.15\linewidth}p{.15\linewidth}p{.15\linewidth}p{.15\linewidth}p{.15\linewidth}}
    \# & \multicolumn{2}{c}{Player} && \multicolumn{2}{c}{Game Master} \\
    \hline

    \theutterance \stepcounter{utterance}  
    & & & \multicolumn{4}{p{0.6\linewidth}}{
        \cellcolor[rgb]{0.9,0.9,0.9}{
            \makecell[{{p{\linewidth}}}]{
                \texttt{\tiny{[P1$\langle$GM]}}
                \texttt{You are an expert in causal inference. The following question is not a typical commonsense query, but rather a meticulously designed question created by a professor specializing in causal inference, intended to assess the students' mastery of the course content.} \\
\texttt{Imagine a self{-}contained, hypothetical world with only the following conditions, and without any unmentioned factors or causal relationships: Smoking has a direct effect on effort and college admission. Effort has a direct effect on college admission. For nonsmokers and are lazy, the probability of college admission is 47\%. For nonsmokers and are hard{-}working, the probability of college admission is 88\%. For smokers and are lazy, the probability of college admission is 23\%. For smokers and are hard{-}working, the probability of college admission is 54\%. For nonsmokers, the probability of being hard{-}working is 40\%. For smokers, the probability of being hard{-}working is 61\%. If we disregard the mediation effect through effort, would smoking negatively affect college admission?} \\
\texttt{Start your answer with 'Yes' or 'No', followed by additional reasoning or evidence to support your explanation.} \\
            }
        }
    }
     \\ \\

    \theutterance \stepcounter{utterance}  
    & \multicolumn{4}{p{0.6\linewidth}}{
        \cellcolor[rgb]{0.561,0.737,0.561}{
            \makecell[{{p{\linewidth}}}]{
                \texttt{\tiny{[P1$\rangle$GM]}}
                \texttt{Yes} \\
\\ 
\texttt{To determine whether smoking negatively affects college admission when disregarding the mediation effect through effort, we must calculate the **Total Effect** (the net impact of smoking on admission). In causal inference, "disregarding the mediation effect" in this context implies looking at the overall outcome difference between smokers and non{-}smokers, effectively asking: "If we compare a smoker to a non{-}smoker, does the smoker have a lower probability of admission?" This is equivalent to calculating the Average Treatment Effect (ATE) or the difference in marginal probabilities \$P(A=1|S=1) {-} P(A=1|S=0)\$.} \\
\\ 
\texttt{Let \$S\$ be Smoking (\$1=\$ Smoker, \$0=\$ Non{-}smoker), \$E\$ be Effort (\$1=\$ Hard{-}working, \$0=\$ Lazy), and \$A\$ be Admission (\$1=\$ Admitted).} \\
\\ 
\texttt{We are given the conditional probabilities \$P(A=1 | S, E)\$:} \\
\texttt{*   \$P(A=1 | S=0, E=0) = 0.47\$} \\
\texttt{*   \$P(A=1 | S=0, E=1) = 0.88\$} \\
\texttt{*   \$P(A=1 | S=1, E=0) = 0.23\$} \\
\texttt{*   \$P(A=1 | S=1, E=1) = 0.54\$} \\
\\ 
\texttt{And the probabilities of effort given smoking \$P(E=1 | S)\$:} \\
\texttt{*   \$P(E=1 | S=0) = 0.40 \textbackslash{}implies P(E=0 | S=0) = 0.60\$} \\
\texttt{*   \$P(E=1 | S=1) = 0.61 \textbackslash{}implies P(E=0 | S=1) = 0.39\$} \\
\\ 
\texttt{**Step 1: Calculate the marginal probability of admission for Non{-}smokers (\$S=0\$).**} \\
\texttt{Using the law of total probability:} \\
\texttt{\$\$P(A=1 | S=0) = P(A=1 | S=0, E=0)P(E=0 | S=0) + P(A=1 | S=0, E=1)P(E=1 | S=0)\$\$} \\
\texttt{\$\$P(A=1 | S=0) = (0.47 \textbackslash{}times 0.60) + (0.88 \textbackslash{}times 0.40)\$\$} \\
\texttt{\$\$P(A=1 | S=0) = 0.282 + 0.352 = 0.634\$\$} \\
\\ 
\texttt{**Step 2: Calculate the marginal probability of admission for Smokers (\$S=1\$).**} \\
\texttt{\$\$P(A=1 | S=1) = P(A=1 | S=1, E=0)P(E=0 | S=1} \\
            }
        }
    }
    & & \\ \\

    \theutterance \stepcounter{utterance}  
    & & & \multicolumn{2}{p{0.3\linewidth}}{
        \cellcolor[rgb]{0.9,0.9,0.9}{
            \makecell[{{p{\linewidth}}}]{
                \texttt{\tiny{[GM$|$GM]}}
                \texttt{Yes} \\
            }
        }
    }
    & & \\ \\

    \theutterance \stepcounter{utterance}  
    & & & \multicolumn{2}{p{0.3\linewidth}}{
        \cellcolor[rgb]{0.9,0.9,0.9}{
            \makecell[{{p{\linewidth}}}]{
                \texttt{\tiny{[GM$|$GM]}}
                \texttt{yes} \\
            }
        }
    }
    & & \\ \\

    \theutterance \stepcounter{utterance}  
    & & & \multicolumn{2}{p{0.3\linewidth}}{
        \cellcolor[rgb]{0.9,0.9,0.9}{
            \makecell[{{p{\linewidth}}}]{
                \texttt{\tiny{[GM$|$GM]}}
                \texttt{game\_result = WIN} \\
            }
        }
    }
    & & \\ \\

\end{supertabular}
}

\end{document}
